% !TeX root = ../../../main.tex
\subsection{Discussion}
\label{ch03:sec:discussion}
Our analyses reveal a robust association between medication non-adherence and the earliest occurrence of any of the three adverse events (mortality, involuntary hospitalization, jail booking) in the short- to medium-term (predicting time until the earliest adverse event starting from 3, 6, or 9 months after the first prescription fill recorded).
Our findings demonstrate that, at each time snapshot, medication non-adherence is associated with earlier adverse event time among patients who have not yet experienced any adverse outcome.
Moreover, at these earlier time frames (3, 6, and 9 months), this effect is particularly pronounced, with non-adherence advancing the onset of adverse outcomes by approximately 3 months on average across our models.
Over the first year, these earlier onsets culminate in an estimated 9\% increase in composite adverse events, underscoring the practical significance of adherence in mitigating severe outcomes.


\ChThreeParagraph{Methodological insights: multiple survival models.}
A central methodological insight from this work is the importance of evaluating multiple survival models in our framework, both for prediction and subsequently in treatment effect estimation.
While Random Survival Forest performed well in our setting, different datasets---especially those with higher dimensional features or more complex event patterns---may favor alternative approaches.
Because survival probabilities and hazard functions form the foundation for counterfactual risk predictions in our framework, careful model selection is indispensable for producing stable Average Treatment Effect (ATE) estimates.

Indeed, our findings show that if a chosen survival model fails to capture the underlying risk dynamics well (as observed with DeepHit in our analysis, where its adverse event prediction performance was noticeably poor), subsequent causal estimates can become unreliable.
Moreover, consistency of ATE estimates across different snapshot times serves as an additional check on a model’s robustness.
In our study, methods with stronger predictive performance (e.g. Random Survival Forest, and Causal Survival Forest) yielded more stable ATE estimates, highlighting a broader analytical insight: reliable causal inferences with right-censored data hinge on trustworthy survival predictions, rendering performance evaluation at multiple time points a critical step in model validation.


\ChThreeParagraph{``Unbundling'' the composite outcome.}
In this study, we treat mortality, involuntary hospitalization, and jail booking as a single composite outcome to leverage well-established causal survival analysis methodologies.
An important next step would be disentangling this composite adverse event and actually modeling the three event types to be distinct.
While competing risks models are widely used in the multiple event type setting, our scenario partially involves events that could happen concurrently at least at the time resolution of the data we have access to (e.g., jail booking and involuntary hospitalization can both happen in the same month, where we do not know which actually happens earlier). 
Also, jail booking and involuntary hospitalization could each happen multiple times for an individual, whereas death of course prevents any future adverse event from happening.
Properly modeling how these three adverse events interact in the presence of censoring would necessitate developing an entirely new analytical framework, as current methods in causal survival analysis do not readily support multiple events like the ones we have (without just bundling them into a larger composite adverse event as we have done).


\ChThreeParagraph{Policy implications and medication type.}
While our subgroup analyses in Section~\ref{ch03:sec:04_results_ite_analysis} and Appendix~\ref{ch03:apd:ite_medication} revealed consistently negative ATE estimates across multiple medication types, we found minimal evidence of meaningful effect heterogeneity in adherence-related outcomes among the top four antipsychotics.
These findings suggest that interventions aiming to enhance medication adherence may not need to target specific medications in isolation.
Instead, policies may achieve greater impact by focusing on the identification and support of high-risk individuals—irrespective of their prescribed medication.
Nevertheless, certain formulations could be inherently more conducive to adherence (e.g., long-acting injections for patients with specific preferences or support systems), so tailoring strategies to individual circumstances remains a fruitful direction for future research.
Ultimately, interventions improving adherence, whether through better communication, usability enhancements, or policy adjustments, may be broadly beneficial, given the uniformity of adherence effects observed across medication types.

From a clinical standpoint, our findings equip healthcare providers with precise, quantitative tools to clearly communicate the specific risks of medication non-adherence to patients.
In our study, given the consistently observed delays in adverse events associated with adherence across different medication types and administration methods, clinicians can leverage this evidence to select medication formulations (injectable versus oral) tailored to individual patient preferences and adherence challenges.
This direct and informed communication can significantly enhance patient awareness and engagement by providing concrete, personalized data that clearly illustrate the potential consequences of non-adherence and benefits of adherence-focused interventions.
In practical terms, clinicians can use these results to better guide treatment discussions, set realistic patient expectations, and design individualized adherence support plans aimed explicitly at incentivizing and improving medication adherence, thus potentially reducing the incidence and severity of adverse outcomes.

From a policy perspective, our results underscore the effectiveness and efficiency of investing in comprehensive adherence-promoting programs rather than narrower, medication-specific interventions.
Policymakers and health administrators can utilize our quantified adherence impacts to prioritize resource allocation towards robust, scalable, and broadly applicable adherence support initiatives.
Examples of such initiatives include structured adherence counseling programs and advanced digital adherence technologies.
These strategically targeted programs can leverage routinely collected administrative and pharmacy data to identify and proactively engage high-risk individuals, such as reaching out to them to check in on their well-being or possibly more directly incentivizing adherence by providing some sort of financial incentive such as a discount or rebate on specific medications.
These targeted problems could significantly mitigate adverse outcomes at the population level and improve overall healthcare quality and efficiency.

\ChThreeParagraph{Actionable interventions informed by adherence quantification.}
Our findings underscore the need for proactive, targeted interventions to enhance adherence.
Effective strategies could include personalized outreach programs, such as scheduled follow-up calls, digital medication reminders, or telehealth sessions designed around patient preferences and needs.
Transitioning suitable patients to long-acting injectable formulations may also mitigate adherence challenges by simplifying medication routines.
Additionally, policy adjustments aimed at streamlining prescription renewals could reduce practical barriers to sustained adherence.
Previous studies have consistently demonstrated that structured adherence interventions significantly lower hospitalization rates and associated healthcare costs \citep{velligan2009expert, kane2016comprehensive}.
Thus, the quantified impact of adherence presented here can directly guide clinicians and policymakers in designing and implementing interventions optimized for maximal clinical benefit and resource efficiency.


\ChThreeParagraph{Other future research directions.}
Although our observational design cannot definitively establish causality, the consistently negative ATE estimates indicate a strong link between non-adherence and earlier adverse events.
Nonetheless, unmeasured factors correlated with adherence may drive this relationship, so these findings do not guarantee that improving adherence will necessarily prevent adverse outcomes.
Given the practical challenges of directly randomizing adherence, a promising avenue is to design RCTs focused on communication-based interventions rather than attempting to force medication use, which would raise ethical concerns.
Such interventions could test whether tailored outreach effectively increase adherence and reduce adverse events.
Future studies should also collect additional clinical and behavioral data, incorporate time-varying confounders, and investigate whether specific subgroups (e.g., by comorbidities or socioeconomic status) are more vulnerable.

While our investigation focuses on individuals with schizophrenia, the general principles extend to other severe mental illnesses that demand sustained medication adherence.
Ultimately, ensuring adherence could delay the onset of life-threatening or socially destabilizing outcomes, reinforcing the importance of ongoing research and intervention development in this critical public health arena.
Although our study is based on data from a single county, the methodological framework we present is broadly applicable.
Future studies could replicate this analysis using data from multiple regions or healthcare systems to assess the consistency of the observed associations across diverse settings.
Additionally, incorporating data from different demographic groups or healthcare infrastructures could provide insights into how contextual factors influence the relationship between medication non-adherence and adverse outcomes.


